\documentclass[../main.tex]{subfiles}
\begin{document}
%==========================================TIÊU ĐỀ TẬP========================================
\begin{table}[h]
    \begin{adjustwidth}{-1cm}{}
        \begin{tabular}{>{\centering\arraybackslash}p{6cm}|>{\raggedright\arraybackslash}p{12.5cm}}
            \multirow{5}{6cm}
            {\\[-40pt]
                \flaregothic\fontsize{20pt}{20pt}\selectfont \centering
                \color{eptype}HISTOIRE\color{black}\\
                \freshbot\fontsize{72pt}{72pt}\selectfont
                \color{epnum}26\\[6pt]
                \cafeteria\fontsize{20pt}{20pt}\selectfont\color{epnum}(D-26)\color{black}
                \\[18pt]
                \flaregothic\fontsize{18pt}{18pt}\selectfont
                \color{epattr}RÉUNION
                \color{black}
            }
             &
            {\fotskip\fontsize{18pt}{18pt}\selectfont \ruby{永}{えい}\ruby{遠}{えん}の\ruby{友}{ゆう}\ruby{情}{じょう}【\ruby{4}{よん}\ruby{期}{き}\ruby{生}{せい}】}
            \\[6pt]
             & {\cafeteria\fontsize{18pt}{18pt}\selectfont You're Our Coco's Four, Forevermore...}                                                                                             \\[4pt]
             & {\vnmsans\fontsize{18pt}{18pt}\selectfont Tình bạn vĩnh cửu II: Gen 4}                                                                                                          \\[8pt]
             & {\caviar{\fontsize{14pt}{14pt}\selectfont Nhân vật: \emp{coco} \emp{kson} \emp{kanata} \emp{watame} \emp{towa} \emp{luna} \emp{kuon2} \emp{kaigou2} \emp{ryuuhou2} \emp{nene}}} \\ & (Dựa theo mẩu truyện tranh ngắn của \href{https://twitter.com/yuu201023}{yuu201023.})
            \\[8pt]
        \end{tabular}
    \end{adjustwidth}
\end{table}
%=========================================TRUYỆN CHÍNH========================================
\color{black}

\txtbx[2]{
    {\color{vol} \uline{Tóm tắt:}} Gần một năm sau ngày chia tay đầy nước mắt trên bờ biển mùa hè năm 2021, thế hệ thứ Tư vẫn không thôi nguôi ngoai nỗi nhớ về người bạn Rồng Kiryu Coco. Một buổi tối ăn ké của Nene tại nhà Kanata đã vô tình làm lộ manh mối về sự tồn tại của một ``ký túc xá đặc biệt'' tại nhà Tổng quản lý Kuon cùng một nhân vật bí ẩn mang tên Kson. Những nét thân thuộc trong lời Nene kể khiến bốn thành viên Gen 4 tìm gặp Kuon để đối chất. Sau những phút giây căng thẳng, nước mắt của Kanata cùng sự can thiệp từ YAGOO đã mở toang cánh cửa hy vọng. Một chuyến đi bất ngờ đến kho hàng của bà Kaien ở ngoại ô Kawachi đã dẫn lối cho cuộc gặp giữa Gen 4 và Kson -- người đã chọn một cái tên mới, nhưng vẫn mang theo ký ức và tình bạn của ngày cũ. Trải qua những lời xác nhận chân thành, các cái ôm nghẹn ngào và một buổi lao động kho bãi hỗn loạn nhưng ngập tràn tiếng cười, Gen 4 hiểu rằng tình bạn của họ không cần phủ nhận con người hiện tại của Kson để tiếp tục bền chặt.
}

\pov{Amane Kanata}

Kiryu Coco-chan.

Một cô nàng Rồng với thân hình cao lớn, bộc trực, nhưng sâu thẳm bên trong lại mang một tâm hồn vô cùng trẻ con.

Cậu ấy rất thích pha trò, phá đám chúng tôi, và cũng là người khởi xướng không biết bao nhiêu vụ rắc rối long trời lở đất trong công ty. Nhiều lúc cậu ấy nghịch ngợm đến mức khiến Kuon-senpai phải phát điên, ôm đầu rên rỉ bất lực.

Thế nhưng, không một ai trong chúng tôi ghét cậu ấy cả. Ngược lại, chúng tôi yêu quý cậu ấy hơn bất kỳ ai.

Cậu ấy là người đã tiếp lửa cho thế hệ chúng tôi ngay từ những ngày đầu bỡ ngỡ bước chân vào nghề. Cậu ấy luôn là nguồn động lực to lớn, không ngừng cổ vũ và kéo chúng tôi đứng dậy mỗi khi vấp ngã. Không chỉ với Gen 4, cậu ấy còn là một tượng đài, một hình mẫu tiên phong mà rất nhiều thành viên khác trong \hll\ luôn kính trọng noi theo.

Chúng tôi đã cùng nhau nếm trải biết bao thăng trầm - từ những giờ phút luyện tập gian khổ, những giọt nước mắt thất vọng cho đến niềm vui vỡ òa trên sân khấu lớn.

Và cả những lần cậu ta thản nhiên văng tục bằng tiếng Anh với vẻ mặt tỉnh bơ nữa. Dù phần lớn thời gian chúng tôi chẳng hiểu cậu ấy đang lẩm bẩm cái gì, ngoại trừ câu cửa miệng quen thuộc vang danh khắp bốn bể: ``Good morning, motherfuckers!''.

Chúng tôi từng ngỡ rằng cả năm đứa sẽ nắm chặt tay nhau đi đến tận cùng giấc mơ.

Rằng chúng tôi sẽ cùng nhau trở thành những VTuber tuyệt vời nhất thế giới.

Thế nhưng\dots\ chỉ vì một tai nạn không đáng có từ những đợt công kích dữ dội của cộng đồng mạng nước ngoài, cậu ấy đã tự nguyện nộp đơn xin rời xa chúng tôi. Cậu ấy rời đi vì không muốn sự hiện diện của mình làm liên lụy và trở thành gánh nặng cho những người ở lại.

Một quyết định dứt khoát đến mức tàn nhẫn, gần như không hề hỏi ý kiến của bốn đứa chúng tôi trước khi công bố.

Tất cả chúng tôi đều bàng hoàng, chết lặng. Đến cả YAGOO-san, chị A-chan hay thậm chí là anh Kuon cũng hoàn toàn không thể lường trước được. Nhưng một khi con Rồng bướng bỉnh ấy đã hạ quyết tâm vì sự an toàn của đồng đội, chúng tôi chẳng còn cách nào khác ngoài việc nghẹn ngào gật đầu chấp thuận.

Và rồi, ngày định mệnh ấy cũng đã đến.

Cuối tháng Sáu, năm 2021.

Tôi vẫn nhớ như in buổi chiều tà ngày hôm ấy - khoảnh khắc mà chúng tôi chính thức nói lời chia tay người bạn Rồng của mình.

\ct

\begin{center}
    \uline{Bãi biển vắng\\Cách thành phố Điện tử ba mươi cây số về phía Bắc.}
\end{center}

Chiều muộn. Gió biển thổi từng cơn rát mặt, mang theo vị mằn mặn của đại dương và tiếng sóng vỗ rì rào vào bờ cát trắng.

``Vậy là\dots\ bà phải đi thật rồi sao, Coco\dots?'' Watame khẽ cất lời, bờ vai nhỏ bé run lên dưới lớp áo choàng.

Coco khoanh hai tay trước ngực, ngước mắt nhìn về phía chân trời đang dần ngả sang màu tím thẫm: ``Ừ. Cũng đến giờ rồi nhỉ.''

Một ngày sau buổi lễ tốt nghiệp chính thức được phát sóng, Coco cùng với anh Kuon đã hẹn bốn đứa chúng tôi ra bãi biển này để nói lời từ biệt riêng tư. Cậu ấy vẫn cố nở một nụ cười toe toét để xua tan bầu không khí ảm đạm: ``Trái Đất này nhỏ bé quá, chật chội không đủ cho một con Rồng như tôi tung hoành đâu nha. Thế nên tôi phải đi tìm một nơi nào đó rộng lớn hơn, đúng chứ?''

``Vậy à\dots'' tôi cúi gằm mặt xuống bờ cát, hai bàn tay siết chặt lấy nhau.

Anh Kuon đứng cạnh đó, thản nhiên chêm vào một câu: ``Chẳng có chỗ nào to hơn cho em sống đâu. Trừ khi em bay thẳng lên sao Mộc thì may ra.''

Chẳng biết anh ấy đang cố tình đùa để xoa dịu hay nói nghiêm túc, nhưng câu bình luận cụt lủn ấy khiến chúng tôi bất giác bật cười trong nghẹn ngào, dù những giọt nước mắt mặn chát đã bắt đầu lăn dài trên khóe mi.

Mặt trời đỏ ối dần chìm xuống đáy biển. Bầu trời tối sầm lại, hệt như cảm giác mất mát đang nuốt chửng tâm can chúng tôi. Làm sao chúng tôi có thể cười nổi trước một cuộc chia ly vĩnh viễn chứ?

Thấy cả bốn đứa cứ gục mặt xuống buồn thiu, Coco liền bước tới, giơ tay đập nhẹ vào vai từng người: ``Ê này! Sao mặt mày ai nấy cứ xị ra như bánh bao chiều thế hả? Tương lai xán lạn vẫn đang chờ đợi chúng ta phía trước cơ mà!''

Tôi nghe thấy tiếng anh Kuon thở dài phía sau: ``Em nói thì hay lắm. Nhưng nhìn các em ấy cúi rạp người xuống như bị gù thế kia, ai mà cười cho nổi?''

``Kuon-paisen thôi vô duyên đi hộ em cái\dots'' Coco lườm anh một cái sắc lẹm. Anh Tổng vội xua tay xin lỗi rồi im bặt. Nàng Rồng thở hắt ra một hơi, giọng chùng xuống: ``Thôi nào\dots\ Chẳng lẽ mọi người không thể vui vẻ tiễn tôi một đoạn đường sao?''

Làm sao chúng tôi vui vẻ cho được? Lời anh Kuon nói tuy có phần thô nhưng lại là sự thật trần trụi.

Chúng tôi ngẩng mặt lên, cố nặn ra một nụ cười gượng gạo nhất trần đời, hai hàng lệ cứ thế tuôn rơi lã chã. Nhưng rồi nụ cười méo xệch ấy cũng vỡ tan thành từng tiếng nấc nghẹn ngào. Chúng tôi không thể chịu đựng thêm được nữa.

Khóc vì sắp phải rời xa người đồng đội chí cốt.

Khóc vì những nỗi niềm uất ức, bất lực bị đè nén suốt bao ngày qua.

Khóc vì đây là lần cuối cùng chúng tôi được đứng bên nhau dưới danh nghĩa Thế hệ thứ Tư trọn vẹn.

``Với lại, biết đâu một cánh cửa mới sẽ mở ra thì sao--''

Không để cậu ấy nói hết câu động viên dang dở, bốn đứa chúng tôi cùng òa khóc, lao thẳng về phía trước như một phản xạ tự nhiên từ tận sâu trong tiềm thức.

Chúng tôi sà tới, ôm chầm lấy người bạn Rồng to lớn bằng tất cả sức bình sinh, vùi mặt vào ngực cậu ấy mà nức nở. Một cái ôm siết chặt đến nghẹt thở, tựa như nếu buông tay ra dù chỉ một giây, người trước mặt sẽ tan biến vào làn gió biển.

Lạ kỳ thay, con Rồng ngang tàng thường ngày hay gạt phắt đi những cử chỉ ủy mị ấy nay lại đứng yên bất động. Cậu ấy vòng hai cánh tay lực lưỡng ôm trọn lấy cả bốn chúng tôi vào lòng.

Tôi ngước mắt nhìn lên qua làn lệ nhòa. Tôi thấy khóe mi của Coco khẽ rung lên. Giọt nước mắt trong suốt lăn dài trên gò má nàng Rồng, dù khóe môi cậu ấy vẫn cố mỉm cười như muốn chối bỏ sự yếu đuối của bản thân: ``Mọi người này\dots\ đừng khóc nữa mà\dots\ Mấy người làm tôi cũng muốn khóc theo rồi đây này\dots\ Đồ ngốc ạ.''

Cứ như vậy, năm người chúng tôi ôm chặt lấy nhau giữa bãi biển lộng gió, lặng lẽ chia sẻ từng hơi ấm cuối cùng. Chúng tôi khóc, khóc mãi, khóc cho đến khi mọi uất ức, tiếc nuối và nghẹn ngào trong lòng tan biến hết vào từng đợt sóng biển hoàng hôn.

\ct

Giờ chia tay cuối cùng cũng điểm.

Khuôn mặt năm đứa chúng tôi đã nhem nhuốc vì nước mắt, nhưng giờ đây trên môi mỗi người đều nở một nụ cười rạng rỡ, đúng như lời Coco dặn.

Anh Kuon bước đến cạnh cậu ấy, nét mặt nghiêm nghị trầm lặng: ``Anh mừng vì các em vẫn luôn yêu thương và gắn kết với nhau như thế này. Nhưng giờ cũng muộn rồi, Coco-chan phải lên đường thôi. Mấy đứa còn điều gì muốn nhắn nhủ không?''

Watame bước lên trước, đôi mắt cừu long lanh: ``Coco-chi, nhớ phải luôn luôn vui vẻ và mạnh khỏe đấy nhé!''

``Được chứ, Chân giò! À, và nhớ cẩn thận đừng để Botan-chan ăn thịt đấy nha!'' Coco nháy mắt tinh quái.

``Tôi có phải đồ ăn đâu chứ!'' nàng Cừu phụng phịu.

Towa siết chặt nắm tay, hất cằm: ``Bà mà dám quên bọn tôi là bà chết với tôi đấy nhé!''

``Bà dọa dẫm mà trông đáng yêu thế kia thì dọa được ai hả? TMT!''

``\textbf{Tôi không phải thiên sứ!}'' cô bạn Ác ma đỏ mặt hét lên.

Luna ngẩng đầu, chu môi dặn dò: ``Bà mà tìm được quán takoyaki nào ngon thì nhớ báo cho tui nha-nanora!''

``Nhất trí luôn! Bữa nào tôi sẽ mời bà ăn tiệc takoyaki ngoài không gian luôn!'' Coco cười lớn.

Anh Kuon quay sang nhìn tôi: ``Còn Kana-chan? Em có điều gì muốn nói không?''

Tôi nhìn sâu vào đôi mắt đỏ pha sắc hồng rực lửa của nàng Rồng. Tôi không cất nên lời. Những gì cần nói, những nỗi niềm sâu kín nhất, chúng tôi đã trao trọn cho nhau qua từng ánh mắt và cái ôm vừa rồi. Coco mỉm cười gật đầu, hiểu thấu tâm can của tôi.

Cậu ấy quay lưng lại với chúng tôi, giơ một bàn tay lên vẫy cao: ``Vậy nha! Kuon-paisen, đi thôi!''

Tôi thoáng ngơ ngác khi thấy anh Kuon bước đi cùng cậu ấy. Thấy vẻ mặt ngờ vực của tôi, Coco quay đầu lại cười khì: ``Anh ấy tiện đường về nhà cùng tôi mà. Đừng có nhìn như thể anh ấy cũng tốt nghiệp theo tôi thế chứ!''

Dứt lời, một luồng ánh sáng rực rỡ bùng lên. Toàn thân Coco biến đổi thành một con rồng phương Tây khổng lồ đầy uy dũng. Cậu ấy vỗ mạnh đôi cánh lớn, tạo nên một cơn lốc xoáy cát bụi rồi vút bay thẳng lên bầu trời đêm lấp lánh muôn ngàn vì sao.

Những chiếc vảy rồng phát sáng rơi lả tả theo làn gió biển.

Tôi đứng chôn chân trên bờ cát, hai tay giơ cao vẫy mãi không ngừng, hét lớn bằng tất cả tình cảm của mình: ``\textbf{Bà đi mạnh khỏe nhé\dots\ Coco!}''

Đó là lần cuối cùng chúng tôi được nhìn thấy bóng hình của cậu ấy.

\ct

Gần một năm đã trôi qua kể từ buổi chiều hoàng hôn trên bờ biển năm ấy.

Sau khi Coco rời đi, bốn người chúng tôi tiếp tục guồng quay công việc tại \hll. Mất đi một trụ cột, chúng tôi học cách trưởng thành nhanh hơn, gắn kết và bao bọc lấy nhau nhiều hơn bao giờ hết.

Tôi được giao trọng trách làm nhóm trưởng của Thế hệ thứ Tư. Tôi trở thành đầu tàu lèo lái mọi hoạt động chung, trở thành điểm tựa vững chãi cho cả nhóm. Dù thi thoảng tôi vẫn bị người hâm mộ và các thành viên khác trôi chọc vì lực tay khỏe như ``Khỉ đột'', hay bị trêu là ``tấm ván phẳng lì'' - nhưng mặc kệ, miễn là tôi có đủ tài năng và năng lượng để dẫn dắt Gen 4 bước tiếp, thì ngoại hình có là cái đinh gỉ gì chứ.

Thế nhưng, dù cuộc sống có bận rộn đến đâu, hình bóng của Kiryu Coco chưa bao giờ phai mờ trong tâm trí chúng tôi.

``\textit{Ước gì bọn mình có thể gặp lại cậu ấy dù chỉ một lần\dots}''

Đó là ước mơ cháy bỏng nhưng dường như là bất khả thi của cả bốn đứa. Chúng tôi ngỡ rằng khoảng cách ấy là vô tận, cho đến một ngày nọ\dots

\ct

\begin{center}
    \uline{Đầu tháng Năm năm 2022.

        Căn hộ của Kanata - Thành phố Điện tử.}
\end{center}

Sau buổi stream chơi game chung cả buổi chiều, bốn đứa chúng tôi chuẩn bị nấu bữa tối. Bất thình lình, Nene mò sang ăn ké. Cô nàng hậu bối tinh nghịch vừa ngồi vào bàn đã mè nheo đòi bón cơm cho Towa, làm con bé Ác ma sợ chết khiếp giãy nảy lên.

Trong lúc rôm rả chuyện trò, Nene bỗng chép miệng đổi chủ đề: ``Mà nhắc đến Kuon-senpai á\dots\ Tháng trước đám NePoLaBo bọn em mới sang nhà anh ấy ăn lẩu đấy!''

``Mấy đứa sang nhà Kuon-senpai chơi à?'' Watame tò mò hỏi.

``Vâng! Chơi game vui lắm! Nhưng đặc biệt nhất là\dots\ bọn em đã gặp lại một người tưởng chừng như không bao giờ còn cơ hội gặp lại nữa!'' Nene hí hửng khoe.

Cả bốn đứa Gen 4 lập tức dỏng tai lên: ``Ai thế?''

``Mano Aloe-chan!''

Cái tên vang lên như một tiếng sét đánh ngang tai. Sau vài giây lục lọi ký ức, Towa reo lên: ``A! Cô bé tiểu quỷ có cặp sừng không cân xứng của Gen 5 đúng không?!''

``Chuẩn luôn ạ! Bây giờ thì cậu ấy đã là Delutaya-chan rồi!'' cô nàng gật đầu lia lịa. ``Và điều sốc nhất là: cậu ấy đang sống ngay tại nhà của Kuon-senpai luôn! Ở tầng hai tòa nhà ấy!''

Bốn đứa chúng tôi đồng thanh há hốc mồm: ``\textbf{Thật thế á?!}''

``Thật một trăm phần trăm! Vì nhà của anh ấy chính là ký túc xá an toàn dành cho các cựu thành viên đã tốt nghiệp mà lị!''

Tôi cảm thấy sống lưng mình lạnh toát. Trực giác mách bảo một điều gì đó vô cùng chấn động. Tôi vội nắm lấy vai Nene, gặng hỏi: ``Nene-chan\dots\ ngoài Aloe-chan ra, ở đó còn có ai sống nữa không?!''

Thấy vẻ mặt nghiêm trọng của tôi, Nene hơi giật mình, ngập ngừng đáp: ``D-Dạ\dots\ hình như ở tầng trên cùng còn có hai người nữa. Một người tên là Mikeneko sống biệt lập với em gái\dots\ còn người kia thì\dots''

``Người kia là ai?!'' tôi sốt sắng dồn ép, khiến cô bé xanh cả mặt.

``Tên là\dots\ Kson ạ.''

Kson? Một cái tên hoàn toàn xa lạ. Nhưng nếu đó là nơi cưu mang những thành viên đã rời công ty, thì ngoài Aloe và chị Rushiara, người thứ ba chỉ có thể là\dots

Watame nhanh nhảu hỏi tiếp: ``Kson là người như thế nào vậy em?''

Nene gãi đầu nhớ lại: ``Ê-tồ\dots\ Xem nào, dáng người rất cao ráo, phong thái ngầu lòi như giang hồ xã hội đen\dots\ bắn tiếng Anh siêu đỉnh, câu cửa miệng là `Good morning, motherfuckers!'\dots\ và đặc biệt là cực kỳ thích sưu tầm đồ lưu niệm có in hình của Kanata-senpai ạ!''

*PHỤT!*

Ngụm nước tôi vừa uống lập tức phun thẳng ra ngoài sàn nhà! Tôi ho sặc sụa, hai mắt mở trừng trừng: ``K\dots\ Khỉ thật\dots\ Không thể nhầm lẫn vào đâu được nữa!''

Cao ráo. Bắn tiếng Anh như gió. Câu chửi thề quen thuộc. Và cuồng merchandise của tôi!

Tôi đập mạnh hai tay xuống bàn, cất giọng run rẩy nhưng chắc như đinh đóng cột trước ba người đồng đội đang nín thở ``Coco-chan\dots\ Cậu ấy đang ở đó!''

Towa, Watame và Luna đồng loạt reo lên kinh ngạc xen lẫn phấn khích tột cùng. Coco vẫn ở đó! Ngay tại căn nhà của Tổng quản lý Kuon.

Chúng tôi quyết định không chần chừ thêm một giây nào nữa. Ngay sáng mai, cả bốn đứa sẽ kéo đến văn phòng để tra hỏi anh Kuon cho ra nhẽ!

\pov{Hikawa Kuon}

\begin{center}
    Sáng hôm sau.

    Văn phòng \hll.
\end{center}

Tám giờ sáng thứ Tư.

Tôi vừa bước vào phòng làm việc, đặt chiếc laptop lên bàn thì đã thấy Ryuuhou đang ngồi vắt vẻo trên chiếc ghế sofa dài ở góc phòng. Con bé đang cầm hộp sữa tươi hút rột rột, tai đeo một bên headphone lướt xem các bản phối âm của ban nhạc underground trên điện thoại.

Tôi nhướng mày, cất tiếng hỏi: ``Ủa Ryuuhou? Nay được nghỉ tiết sớm hay sao mà mò lên tận công ty anh ngồi ăn sáng thế kia?''

Con bé ngẩng lên, tháo tai nghe xuống toe toét cười: ``Sáng nay trường em khối trên thi tốt nghiệp nên bọn em được nghỉ tiết đầu anh ạ! Em tạt qua ăn ké cái bánh mì với hộp sữa, tiện thể đợi lát nữa anh rảnh chở em sang phòng tập nhạc bên Ochanomizu luôn.''

``Lại coi anh mày như tài xế xe ôm miễn phí rồi đấy,'' tôi lắc đầu bật cười, kéo ghế ngồi xuống. ``Bên đó chuẩn bị diễn show mới à?''

``Vâng, bọn em đang khớp lại mấy bài rock mới! Mà dạo này công việc ở văn phòng thế nào rồi Onii? Hôm trước em nghe Kaigou-nee bảo chi nhánh EN chuẩn bị ra mắt thế hệ mới bận rộn lắm hả anh?''

``Ừ, việc ngập đầu ra đây này\dots''

Hai anh em chúng tôi còn chưa kịp nói hết câu thì cánh cửa văn phòng bất ngờ bị đẩy mạnh ra một tiếng *RẦM* chói tai.

Bốn thành viên của Thế hệ thứ Tư - Kanata, Towa, Watame và Luna - bước vào phòng theo một hàng ngang. Trên gương mặt họ toát ra một thứ sát khí ngùn ngụt đến mức khiến nhiệt độ trong phòng như tụt xuống cả chục độ.

Ryuuhou giật mình trố mắt nhìn ra cửa. Thấy bộ dạng đằng đằng sát khí của bốn cô nàng idol, con bé vội đặt hộp sữa xuống bàn, nhanh chóng đứng dậy bước lại gần bàn làm việc của tôi.

Kanata chống hai tay lên mép bàn tôi, người hơi nhoài về phía trước, giọng gằn từng chữ: ``Kuon-senpai, chúng ta cần nói chuyện nghiêm túc.''

Tôi chớp mắt, cố giữ vẻ điềm tĩnh thường ngày: ``Có chuyện gì mà mấy đứa căng thẳng thế?''

Towa không thèm vòng vo, khoanh tay cắt ngang bằng một câu hỏi đanh thép: ``Coco-chan đang ở nhà anh đúng không?!''

Tim tôi dường như có chút giật mình khi câu hỏi đó rơi vào không khí căng thẳng của văn phòng, nhưng tôi cố gắng kiềm chế mọi cảm xúc, giữ cho vẻ mặt vẫn thản nhiên, không lộ ra bất cứ dấu vết nào. Ban đầu, tôi chỉ nghĩ rằng có lẽ họ tình cờ nghe được một lời đồn thổi vô tình nào đó từ đồng nghiệp trong công ty, hay chỉ là những phỏng đoán vu vơ từ cộng đồng người hâm mộ bên ngoài. Nhưng nhìn bốn cô gái đứng trước mặt, gương mặt đầy vẻ nghiêm trọng đến mức dường như có thể thiêu đốt mọi thứ, tôi bỗng dưng nghi ngờ phỏng đoán ban đầu của mình. Tại sao họ lại trông dồn dập và quyết liệt đến thế?

Chưa dừng lại ở đó, đằng sau lớp vẻ mặt cương nghị ấy, tôi còn đọc thấy được nỗi tuyệt vọng đang âm ỉ cháy trong mắt từng người. Họ không chỉ đơn thuần là tò mò hay nghi ngờ, mà họ đã đặt tất cả niềm tin và hy vọng vào câu trả lời của tôi, như thể đây là cơ hội cuối cùng để họ lấy lại điều gì đó quan trọng đã mất.

Chú YAGOO đã từng nói với tôi vài tuần trước rằng sớm muộn gì bốn đứa Gen 4 cũng sẽ tìm đến, và lời tiên tri đó quả nhiên đã ứng nghiệm không sai một ly nào. Đợt hội ngộ tiếp theo của dự án ``Như chưa hề có cuộc chia ly'' mà chúng tôi đã chuẩn bị từ lâu, hóa ra lại không cần chúng tôi chủ động mời gọi, mà chính những người tham gia đã tự tìm đến cửa văn phòng của tôi.

Và đúng như những gì tôi đã đoán trước, chỉ vài giây sau khi tôi cố gắng né tránh câu hỏi, Watame-chan đã nhanh nhảu lên tiếng xác nhận: ``Nene-chan có nói với chúng em như vậy.'' Tất nhiên, với những chứng cứ mà bốn đứa nắm trong tay thì vẫn chưa đủ để đưa ra bất kỳ kết luận chính thức nào, và thế là cuộc đối đáp căng thẳng giữa tôi và bốn cô gái bắt đầu diễn ra.

Luna nhảy cẫng lên một chút, đôi tay siết chặt vào nhau, giọng nói hốt hoảng xen lẫn sự van nài: ``Anh nói thựt đi, đừng giấu giếm gì nữa, khai ra hết tất cả nhữn zì anh bí cho bọn em-nora!''

Watame gật đầu lia lịa theo, người hơi nghiêng về phía bàn làm việc của tôi, giọng run rẩy nhưng đầy quyết tâm: ``Anh không thể giữ bí mật này mãi được đâu, bọn em đã biết hết rồi!''

Tôi giơ hai tay lên đầu, lắc đầu ngao ngán trước cảnh tượng bốn cô gái đứng thành hàng ngang trước mặt với gương mặt đầy sát khí: ``Trời ơi, bốn đứa đứng đây trông cứ như băng đảng xã hội đen đòi nợ thuê vậy\dots''

Kanata càng đập mạnh bàn thêm một lần nữa, giọng hét lớn như muốn đục xuyên cả bức tường văn phòng: ``Khai ra hết mau đi, Kuon-senpai! Đừng có cố mà che giấu nữa!''

Ryuuhou đứng ngay cạnh tôi, khoanh hai tay trước ngực. Nhận thức được tính chất bảo mật tối quan trọng của khu ký túc xá và sự an toàn của những người đang sống tại đó, con bé lập tức phối hợp cùng tôi, cất giọng chất vấn ngược lại: ``Khoan đã nào các chị\dots\ Tự dưng sáng sớm các chị xông vào với sát khí ngùn ngụt thế làm bọn em giật cả mình. Chuyện của Coco-san năm ngoái ai chẳng buồn và tiếc nuối, nhưng cớ sao các chị lại đùng đùng chạy đến hỏi Onii câu đó?''

Tôi gật đầu phụ họa em gái, ngả người ra sau ghế nhìn xoáy vào bốn cô gái: ``Ryuuhou nói đúng đấy. Ai nói cho mấy đứa biết chuyện đó?''

``Nene-chan đã kể hết cho bọn em nghe rồi!'' Watame đáp.

Ryuuhou nheo mắt, ra vẻ hoài nghi đầy sắc sảo: ``Nene-san kể á? Chị ấy thì chúa bày trò quậy phá với thêu dệt chuyện trên trời dưới đất rồi. Chị ấy đã kể những gì mà làm các chị tin sái cổ vậy?''

Tôi tiếp lời em gái để kiểm tra mức độ xác thực của họ: ``Đúng thế, em ấy kể những gì?''

Luna vừa nhảy cẫng lên vừa chỉ vào mặt tôi, giọng run rẩy: ``Nene-chan kể với bọn em rằng em í đã từng đến nhà của anh chơi rồi! Mấy chuyện đó là thật đúng không ạ?''

Tôi khoanh tay, mặt không biểu cảm gật đầu: ``Ừ, đúng. Nene-chan có đến thăm nhà anh vài lần. Tiếp đi.''

Nàng Cừu cũng vội vàng chen ngang, đôi mắt cừu long lanh đầy hy vọng: ``Rồi tại đó em ấy còn gặp lại Mano Aloe-chan, cô bạn cũ của Thế hệ thứ Năm nữa! Đúng không anh?''

``Tiếp đi, có gì nữa thì nói hết ra.''

Towa bước tới cạnh Kanata, tay siết chặt đến trắng bệch: ``Em ấy cũng nói rằng ngoài Aloe-chan ra, ở đó còn có hai người khác đang sống cùng. Mấy chuyện này đều là thật đúng không?''

Trong lòng tôi thầm thở phào: không sai một ly nào, đúng hết cả. Nhưng tôi vẫn cố giữ bộ mặt lạnh tanh, muốn tiếp tục tra khảo để xem họ có thật sự nắm rõ mọi chuyện hay chỉ nghe lóm được vài câu đồn thổi vu vơ. Thành thật mà nói, ai mà chẳng có thể nghe ai đó kể rồi phóng đại lên, chuyện lớn hóa nhỏ, nhỏ hóa to. Tôi cần phải thấy đủ bằng chứng mới chịu tin, không phải vì khó tính, mà vì mọi thứ liên quan đến những cựu thành viên đều phức tạp hơn họ tưởng rất nhiều.

Tôi gõ ngón tay lên mặt bàn, hỏi như thể đang thẩm vấn: ``Hai người đó tên gì? Nói cho anh nghe.''

Towa ngẩng cao đầu, trả lời liền một mạch không chần chừ: ``Một người tên là Mikeneko, còn người kia tên là Kson! Nene-chan nhớ rõ tên cả hai!''

Ryuuhou gật đầu ra vẻ hài lòng, nhưng vẫn không thả lỏng thái độ: ``Đúng rồi, có hai người đó thật. Còn gì nữa không ạ? Nene-san miêu tả họ là những người như thế nào ạ?''

Watame gãi đầu nhắc lại từng chữ Nene đã kể, giọng hơi ngập ngừng: ``Ê-tồ\dots\ Nene-chan có tả người tên Kson đó là\dots\ là một người rất cao lớn, dáng vẻ ngầu lòi, tính cách cũng khá ngang tàng, bặm trợn lắm\dots''

Luna chớp mắt nghiêng đầu, quên mất hết thông tin Nene chia sẻ: ``Tôi đâu có nhớ em ấy nói đến tính cách của người đó đâu nhỉ-nora? Tôi chỉ nhớ có nghe tên Kson thôi!''

Cô bạn Cừu quay sang nhìn bạn cùng nhóm, mặt đầy ngạc nhiên: ``Hể? Cái gì mà quên nhanh thế? Tôi tưởng kiểu người xã hội đen nào mà chả có tính cách ngang tàng, bặm trợn thế kia chứ!''

Tôi bật cười khè khè, gật đầu công nhận: ``Em ấy tả đúng đấy, Kson đúng là kiểu người như vậy. Còn gì nữa không?''

Towa vội vàng bổ sung thêm, sợ bỏ sót thông tin nào: ``Nene-chan còn nói người đó giỏi tiếng Anh siêu đỉnh, nói chuyện với ai cũng xài tiếng Anh lưu loát lắm!''

Ryuuhou nhún vai, vẫn giữ vẻ mặt thản nhiên: ``Nhà Hikawa bọn em có nhiều người giỏi tiếng Anh mà, Kson-san chỉ là một trong số đó thôi. Có thông tin gì đặc biệt hơn không ạ?''

Kanata đột nhiên im bặt, môi mím chặt, cả người đứng yên như hóa đá. Cả phòng ai cũng nhìn về phía cô.

Từ góc bàn, tôi bất chợt thấy gò má của ``Thiên sứ lực lưỡng'' bỗng dưng đỏ bừng lên như vừa bị nướng lửa. Tại sao lại thế nhỉ?

Rồi chỉ vài giây sau, tôi liền hiểu ra lý do. Em ấy ngẩng mặt lên, giọng lắp bắp đến mức nói không nên lời: ``Em ấy\dots\ Nene-chan còn nói rằng\dots\ n-người Kson đó\dots\ r-rất thích sưu tầm tất cả đồ vật có in hình\dots\ c-của em ạ.''

Thật không ngờ, cái chi tiết đó cũng đúng y chang. Kson mới chỉ bắt đầu đặt hàng lưu niệm của Kanata mấy tuần gần đây thôi, và em ấy chẳng bao giờ ngại chia sẻ với mấy người ở nhà về niềm đam mê sưu tầm đó. Dù đã biết trước, nhưng lúc này nghe Kanata nói ra, tôi vẫn không khỏi ngạc nhiên. Mọi đặc điểm mà bốn người kể ra, đều trùng khớp hoàn toàn với Kson, cũng chính là Kiryu Coco, người bạn Rồng mà chúng tôi đã giấu kín gần một năm qua.

Tôi khẽ gật đầu: ``Được rồi,'' rồi tiếp tục hỏi: ``Ngoài Nene-chan ra mấy đứa còn nghe được ai nói gì không?''

Nhóm Gen 4 đều lắc đầu đáp lại: ``Không ạ.''

Nghe đến đó, tôi và Ryuuhou khẽ liếc nhìn nhau trong tích tắc. Dù trong lòng cả hai đều biết tỏng lời Nene kể là sự thật trăm phần trăm, nhưng ngoài mặt chúng tôi vẫn phải giữ một vẻ lạnh tanh.

Em tôi lắc đầu bẻ lại ngay tắp lự: ``Chỉ có thế thôi ạ? Trùng hợp thôi các chị ơi! Ở ngoài kia người mê phong cách yakuza với bắn tiếng Anh đầy ra đấy. Với lại, hàng lưu niệm của Kanata-san bán khắp nơi, ai là fan chẳng mua được vài món sưu tầm? Chẳng lẽ cứ ai dáng cao, thích chửi thề tiếng Anh với cuồng đồ của chị ấy thì đều là Coco-san hết sao?''

``C-Chị ngĩ thế, nora\dots?'' Luna gãi đầu, lắp bắp trả lời.

Tôi gật đầu công nhận lập luận của em gái, rồi nhìn thẳng vào bốn người họ, buông một câu dội gáo nước lạnh: ``Ryuuhou nói không sai đâu. Chưa đủ chứng cứ để kết luận đâu.''

``Cái gì cơ ạ?!'' cả bốn cô gái sững sờ kêu lên.

``Nene là một đứa thích trêu chọc và hay thêm mắm dặm muối,'' tôi bình thản giải thích. ``Giống nhau về tính cách không có nghĩa Kson và Coco là cùng một người. Ngoài những lời kể vu vơ của Nene-chan ra, các em còn có bằng chứng xác thực nào khác để khẳng định Coco-chan đang ở nhà anh không?''

Câu hỏi hóc búa của tôi và Ryuuhou như một nhát búa giáng thẳng vào niềm tin của họ. Cả căn phòng chìm vào sự im lặng nghẹt thở. Họ hoàn toàn tắc tịt, không thể đưa ra thêm bất kỳ chứng cứ thực tế nào.

Xem chừng ca này với họ sẽ rất khó đây. Thế bàn cờ đã bị lật ngược rồi.

\pov{Amane Kanata}

Khỉ thật! Khỉ thật đấy!

Anh ấy và cả Ryuuhou-chan nói không sai một lời nào. Chúng tôi chỉ nghe qua lời kể của Nene, ngoài ra chẳng có bất kỳ tấm ảnh hay bằng chứng nào trong tay cả. Nhỡ đâu Nene chỉ đang bày trò trêu chọc chúng tôi thì sao?

Cảm giác hy vọng tột cùng vừa lóe lên đã bị dập tắt không thương tiếc.

Nỗi tuyệt vọng nuốt chửng lấy tâm trí tôi. Lồng ngực tôi nghẹn ứ, hai hàng nước mắt trào ra không thể kiểm soát.

Cảm giác ngỡ như đã chạm tay tới cậu ấy\dots\ vậy mà giờ đây lại vụt bay mất!

``Không chịu đâu\dots! \textbf{Tôi không chịu đâu mà!}''

Tôi hét lên trong cơn quẫn trí, hai nắm đấm thép siết chặt lại rồi nện bôm bốp xuống mặt bàn làm việc! Chiếc bàn gỗ dày rung chuyển bần bật, mặt bàn nứt toác ra làm đôi dưới lực đấm khủng khiếp của tôi. Giấy tờ, bút viết văng tung tóe khắp sàn nhà.

``Ối mẹ ơi!''

Ryuuhou-chan đứng ngay cạnh hoảng hồn nhảy lùi phắt lại ba bước, hai mắt tròn xoe nhìn chằm chằm vào mặt bàn nứt toác, miệng há hốc kinh hãi: ``Mặt bàn gỗ sồi dày cộp thế kia mà nứt làm đôi\dots\ Lực tay năm mươi mốt ký của Kanata-san ghê thật đấy anh Kuon\dots''

Anh Kuon hoảng hốt chộp vội lấy chiếc laptop kéo về phía mình, trố mắt nhìn tôi như nhìn thấy quái vật, mồ hôi hột túa ra đầy trán, miệng thốt lên: ``Em ấy lúc bị dồn ép vào đường cùng trông đáng sợ thật\dots!''

Towa và Watame vội vàng chạy tới giữ chặt lấy hai tay tôi, trong khi Luna hoảng sợ ôm lấy vai tôi dỗ dành.

Tôi gục đầu xuống mảnh bàn vỡ, òa khóc nức nở như một đứa trẻ bị cướp mất món đồ chơi quý giá nhất trần đời: ``Tại sao chứ\dots\ Bọn em chỉ muốn gặp lại cậu ấy thôi mà\dots\ Coco ơi\dots''

\pov{Hikawa Kuon}

Nhìn Kanata gục đầu trên mảnh bàn vỡ, tiếng khóc nghẹn ngào vang lên như xé toạc cả không khí trong phòng, trái tim tôi bỗng chùng xuống một nỗi đau xót khó tả. Tình bạn mà bốn cô gái Gen 4 dành cho nhau, cái tình cảm mà họ dành cho Coco quả thực là một điều gì đó vô cùng sâu sắc, mạnh mẽ đến mức không thể diễn tả bằng lời. Tôi thừa nhận, trong suốt những năm làm quản lý tại đây, tôi chưa từng chứng kiến một mối liên kết nào bền chặt và cháy bỏng đến thế.

Tôi hiểu rõ rằng, nếu không có cách nào để họ gặp lại Coco, thì Kanata này, cùng với ba người đồng đội còn lại, sẽ tìm mọi cách, dồn tất cả sức lực để làm nên điều mà chính họ cũng từng nghĩ là không thể. Từ ngày đầu tiên cả hai cùng gia nhập \hll, tôi đã nhận ra Kanata và Coco chính là một cặp bài trùng tuyệt hảo, họ bổ trợ cho nhau, cùng nhau vượt qua biết bao thử thách, gắn bó như hình với bóng. Sức mạnh của tình bạn giữa họ không phải là thứ có thể bị ngăn cách bởi bất kỳ quy định hay khoảng cách nào.

Bản thân tôi cũng khao khát được thấy họ đoàn tụ, được nhìn thấy nụ cười rạng rỡ của Thế hệ thứ Tư trọn vẹn một lần nữa. Bởi vì không lâu trước đây, nhóm thế hệ thứ Năm đã có một cuộc hội ngộ đầy thành công, mọi thứ diễn ra vô cùng suôn sẻ, để lại những kỷ niệm đẹp cho tất cả mọi người. Tôi cũng mong rằng Gen 4 sẽ có được cơ hội tương tự, để xua tan đi nỗi nhớ nhung đã dày vò họ suốt gần một năm qua.

Nhưng quy tắc công ty vẫn là quy tắc, những văn bản và quy định bảo mật mà tôi đã đặt ra khi nhận vị trí này không cho phép tôi tùy tiện làm trái. Liệu rằng tôi có quá cứng nhắc, quá rập khuôn với những quy định đó, không biết linh hoạt khi đứng trước tình cảm chân thành của các em ấy?

Thực sự thì, cái chuyện nhóm NePoLaBo tình cờ gặp lại Delutay tại nhà tôi cũng hoàn toàn là một sự ngẫu nhiên, không hề nằm trong kế hoạch hay dự tính nào của tôi. Ngay sau sự việc đó, sếp của tôi - chú Tanigo - đã gọi điện cho tôi, và tôi đã kể lại mọi thứ một cách trung thực, không giấu giếm một chi tiết nào. Từ trước đến nay, tôi vốn là một người không biết cách nói dối, mọi việc xảy ra đều phải trình bày rõ ràng với cấp trên.

Nhưng với trường hợp của Gen 4, mọi chuyện lại hoàn toàn khác phức tạp. Nếu tôi mở cửa cho họ gặp lại Coco, liệu rằng những người liên quan khác có thật sự chấp nhận và ủng hộ quyết định này không? Những hệ quả, những vấn đề có thể phát sinh sau đó sẽ còn khó lường hơn rất nhiều, tôi không thể lơ là, không thể tùy tiện đưa ra quyết định khi chưa cân nhắc kỹ lưỡng mọi mặt.

Ryuuhou đứng cạnh thấy cảnh tượng ấy cũng lặng người đi. Nhìn thấy những giọt nước mắt lăn dài và bờ vai run rẩy của Kanata, con bé không thể tiếp tục giữ vẻ mặt nghiêm nghị thăm dò được nữa.

Em tiến lại gần, khẽ đặt tay lên vai Kanata vỗ về an ủi: ``Kanata-san\dots\ chị đừng khóc nữa mà\dots''

Đoạn, con bé quay sang nhìn tôi, ánh mắt khẽ lay động như muốn nhắn nhủ: ``Họ thật lòng đến mức này rồi đấy, Onii\dots''

Towa bước lên trước bàn tôi, ánh mắt rưng rưng nhưng tràn đầy sự kiên định: ``Kuon-senpai\dots\ em biết bọn em không có bằng chứng vật lý. Nhưng lúc Nene-chan kể về cuộc hội ngộ với Aloe-chan, em ấy đã cười rất tươi và rơm rớm nước mắt. Nene-chan có thể hay đùa nhây, nhưng em ấy là một người sống rất tình cảm. Em ấy tuyệt đối không bao giờ lấy sự mất mát và nỗi nhớ đồng đội của bọn em ra làm trò đùa đâu! Những gì em ấy thấy là thật, và niềm tin của bọn em dành cho Coco cũng là thật!''

Watame và Luna cũng gật đầu lia lịa, ánh mắt như van nài một sự thấu hiểu từ tôi.

Cô nàng Công chúa xứ Kẹo vừa giậm nhẹ chân vừa đưa tay chống cằm, giọng nói đầy chắc chắn: ``Nó khum hề có nụ cười nham hiểm gì đâu-nanora! Em thấy Nene-chan nói chuyện nghim túc lắm á.''

Tôi nhướn mày ngạc nhiên, hơi nghiêng người về phía trước: ``Thế á?'' Với bản tính hay đùa nghịch và thường xuyên tung những trò đùa tinh quái của em ấy, điều này thực sự khiến tôi bất ngờ. Tôi vẫy tay thúc giục: ``Tiếp đi, còn gì nữa không?''

Nàng Cừu, người vẫn đang vỗ nhẹ lưng để an ủi Kanata đang thút thít, nuốt nước bồi rồi tiếp lời với giọng run run: ``Lúc sau bữa ăn lẩu hôm đó, Nene-chan còn kể lại từng câu chuyện nhỏ nhặt mà em ấy đã trải qua cùng Delu-chan ở nhà Kuon-senpai. Cứ từng chi tiết một, nghe cái nào cũng chân thực như em ấy đang xem lại một đoạn phim vậy. Nghe đến đâu bọn em trong nhóm NePoLaBo cũng cảm thấy từng chuyện đó đều là thật.''

Luna gật đầu lia lịa đồng tình, hai má hồng lên vì xúc động: ``Thật sự luôn đó anh! Nghe Nene-chan kể, bọn em cũng không hề cảm thấy em ấy đang nói dối hay bịa chuyện đâu-nora. Nước mắt em ấy còn rơm rớm lúc kể về lúc gặp lại Aloe-chan nữa!''

Phải công nhận rằng Nene đúng là một cô nàng tinh nghịch không thể chê được, nhưng bên trong lớp vỏ đùa giỡn đó, em ấy lại luôn sống hết mình với những tình cảm chân thành của bản thân. Ngay từ khi Delutaya gặp phải sóng gió trong công việc và phải rời đi, em ấy đã là người trong nhóm NePoLaBu lo lắng nhất, cũng là người buồn rầu nhất suốt bao tháng ngày qua. Nỗi nhớ về người bạn cũ đã từng gắn bó cùng nhau stream, cùng nhau vượt qua bao khó khăn, đó là điều mà em ấy hiểu rõ hơn ai hết. Nếu có ngày nào em ấy dám đem câu chuyện ``tình cờ gặp lại Aloe'' ra làm trò đùa để trêu chọc người khác, chắc chắn người đầu tiên cảm thấy day dứt và buồn bã sau đó sẽ chính là Nene. Với em ấy, nỗi nhớ về những người bạn đồng đội là điều thiêng liêng nhất, không đời nào em ấy có thể đem nó ra để đùa cợt, để làm trò tiêu khiển cho vui miệng được.

Cho nên tất cả những gì em ấy đã kể với bốn cô gái hôm qua, từ việc gặp lại cô bạn ở nhà tôi, cho đến những chi tiết nhỏ nhặt về Kson và Mikeneko, đều là những gì em ấy được chứng kiến tận mắt, chứa đựng đầy đủ những xúc cảm chân thật nhất. Và đó cũng là lý do tại sao Luna, Towa, Watame và Kanata lại tin tưởng tuyệt đối vào lời em ấy, không có chút nghi ngờ nào.

Tôi từ từ ngước mắt nhìn bốn người đang đứng trước mặt mình.

Luna với đôi mắt to tròn vẫn đang nhìn tôi chằm chằm, Towa cũng không thua kém với ánh mắt van nài, Watame thì gương mặt còn rưng rưng những nỗ lực kìm nén nước mắt.

Và rồi đến Kanata, người vừa mới nứt cả một mặt bàn gỗ vì sức mạnh phi thường, giờ đây gục đầu trên mảnh bàn vỡ, vai vẫn còn run rẩy từng hồi. Chính xác hơn, em ấy là người khao khát được gặp lại Coco nhất trong bốn người. Giữa cô nàng thiên sứ và cô nàng rồng đó có một sự gắn kết không thể tả nổi, một sự kết nối đã được tôi chứng kiến từ ngày đầu tiên hai người cùng gia nhập công ty, cùng nhau vượt qua bao thử thách để trưởng thành như ngày hôm nay.

\dots Tự dưng tôi cảm thấy mình thật tệ. Như thể tôi đang là người tước đi cơ hội duy nhất để họ được đoàn tụ, để bốn người một nhà của \textit{holoForce} được đứng chung dưới một mái nhà sau gần một năm chia xa.

Nhưng tôi có thể làm gì được chứ? Quy định bảo mật của công ty, những điều khoản về sự an toàn của các cựu thành viên đang được tôi bảo vệ tại nhà mình, tất cả những gì tôi phải tuân theo không cho phép tôi tùy tiện mở cửa cho họ gặp mặt. Chú YAGOO cũng chưa có bất kỳ thông báo hay động thái chính thức nào sau cuộc hội ngộ thành công của nhóm NePoLaBo vừa qua, tôi đâu thể tự ý quyết định mọi thứ được?

Trời ơi, sao mà tình thế này khó xử đến thế! Đứng giữa tình cảm chân thành của bốn cô gái và trách nhiệm của một người quản lý, tôi thực sự không biết phải xoay sở như thế nào mới phải.

Tôi đứng trước thế tiến thoái lưỡng nan. Luật lệ và quy tắc bảo mật của công ty buộc tôi phải giữ khoảng cách, nhưng tình cảm chân thành của bốn cô gái trước mặt lại khiến tôi không nỡ quay lưng.

Đúng lúc ấy, chiếc điện thoại trên bàn tôi bỗng rung lên bần bật. Trên màn hình hiện lên cuộc gọi video từ ``YAGOO''.

Tôi vội cầm máy bước ra mép cửa sổ, bấm nhận: ``Alo, Hikawa Kuon nghe đây ạ.''

Khuôn mặt tươi cười của sếp tôi hiện lên trên màn hình: ``Chào Kuon-kun. Tình hình văn phòng sáng nay thế nào-- Ồ, đằng sau cậu có chuyện gì mà ồn ào thế? Ai đang khóc vậy?''

Tôi hướng camera về phía bàn làm việc, nơi Kanata vẫn đang thút thít trong vòng tay của ba người bạn và sự an ủi của Ryuuhou. Tôi thở dài, tóm tắt toàn bộ câu chuyện vừa diễn ra.

Nghe xong, chú Tanigo im lặng vài giây, rồi khẽ mỉm cười ấm áp qua màn hình: ``Tôi hiểu rồi. Nhóm Gen 5 đã có một khởi đầu rất tuyệt vời. Và tôi nghĩ, đối với Gen 4, thời điểm thích hợp cũng đã đến rồi đấy. Cậu không cần phải tiếp tục làm kẻ phản diện giấu giếm bọn họ nữa đâu.''

Tôi ngạc nhiên: ``Ý chú là\dots?''

``Cho các cô gái gặp lại Kiryu-san hay không là toàn quyền quyết định ở cậu, Kuon-kun. Tôi hoàn toàn ủng hộ!''

Nói xong, ông ấy nháy mắt một cái rồi cúp máy, để lại tôi đứng ngơ ngác giữa căn phòng.

Tôi quay đầu lại. Towa, Watame và cả Ryuuhou đang nín thở nhìn tôi chờ đợi phán quyết cuối cùng. Tôi thở dài một hơi thật sâu, khóe môi khẽ cong lên một nụ cười nhẹ nhõm: ``Được rồi. Nếu sếp đã bật đèn xanh\dots\ thì anh chịu thua mấy đứa.''

\pov{Amane Kanata}

Tôi từ từ ngẩng khuôn mặt nhem nhuốc nước mắt lên.

Anh Kuon bước tới trước mặt bốn đứa chúng tôi, hít một hơi thật sâu rồi dõng dạc tuyên bố: ``Thế hệ thứ Tư\dots\ bốn đứa các em sẽ là những người tiếp theo chính thức tham gia chương trình `Như chưa hề có cuộc chia ly'!''

Tôi chết lặng. Toàn thân tôi run bắn lên: ``Th\dots\ Thật thế sao anh?!''

Towa và Watame reo hò ầm ĩ, Luna nhảy cẫng lên ôm lấy cổ tôi. Ryuuhou đứng bên cạnh cũng nở nụ cười rạng rỡ, vỗ tay chung vui: ``Tuyệt vời quá! Chúc mừng các chị nhé!''

Tôi không kìm được xúc động, lao thẳng tới ôm chầm lấy anh Kuon, miệng rối rít cảm ơn anh trong nước mắt: ``Cảm ơn anh! Cảm ơn anh nhiều lắm Kuon-senpai!''

Anh Tổng ngượng ngùng gãi đầu: ``Thôi nào, anh có làm được gì đâu, ban nãy còn làm em khóc cơ mà\dots''

Đúng lúc đó, điện thoại của anh Kuon rung lên một tiếng *TINH*. Anh mở ra xem rồi bật cười lớn: ``Nhắc tào tháo là tào tháo tới liền. Kson vừa nhắn tin cho anh, bảo chiều nay cô ấy sẽ xuống kho hàng của mẹ anh ở ngoại ô để phụ bốc vác.''

Cả bốn đứa chúng tôi tròn xoe mắt: ``Chiều nay luôn sao ạ?!''

``Ừ. Nếu mấy đứa muốn gặp em ấy sớm, chiều nay theo anh xuống kho phụ một tay luôn. Chịu cực một chút nhé!''

Ryuuhou tiếc nuối tặc lưỡi: ``Tiếc quá, chiều nay em có lịch tập dợt với ban nhạc rồi nên không đi hóng vui cùng mọi người được. Kanata-san với các chị nhớ bắt đền chị Kson một bữa ra trò đấy nhé!''

``Chắc chắn rồi em!'' tôi gật đầu thật mạnh, rồi cả bốn đứa đồng thanh hô vang đầy hào khí: ``Bọn em đi! Đi liền chiều nay luôn ạ!''

Kson! Bà hãy đợi đấy! Bọn tôi đến gặp bà đây!

\pov{Kson}

Trời ơi, chẳng thể nào tin nổi vào tai mình được! Thật sự là bọn họ sắp đến gặp tôi sao? Những người bạn đồng hành thân thiết nhất của tôi từ ngày đầu tiên cùng bước vào công ty đây mà!

Ngay cả cô nàng cừu non Watame hay nàng công chúa kẹo ngọt Luna, cả Towa với biệt danh TMT nữa\dots\ và đương nhiên không thể thiếu được cô bạn thiên sứ lực lưỡng - khỉ đột Kanatan của tôi rồi!

Đã bao lâu rồi nhỉ? Từ ngày tôi rời đi, chúng tôi chưa từng có cơ hội gặp mặt trực tiếp như thế này. Tôi cứ tự hỏi không biết trong suốt thời gian xa cách ấy, từng người trong họ đã trưởng thành và thay đổi ra sao nhỉ?

Suy nghĩ một chút, tôi bỗng nảy ra một ý định thú vị. Không phải để thử xem họ có nhận ra ``Coco'' dưới một lớp vỏ mới hay không -- Kson không phải lớp vỏ, mà là cái tên tôi đã chọn để sống tiếp. Nhưng tôi cũng không cần vờ như những năm tháng cũ chưa từng tồn tại. Có lẽ chỉ cần để họ gặp tôi đúng như bây giờ: Kson, với những ký ức chung và vài thói quen mà tôi vẫn muốn giữ lại.

He he\dots\ cứ thế đi, tôi sẽ đợi mọi người ở kho hàng. Nếu họ nhận ra điều gì thân thuộc, thì đó sẽ là vì tình bạn của chúng tôi chưa từng bị một cái tên xóa đi.

\pov{Hikawa Kuon}

Chiều hôm đó.

Kho hàng công ty vật liệu xây dựng của mẹ tôi - nằm ở vùng ngoại ô xã \ruby{Ankei}{An Khánh}, cách trung tâm thành phố Kawachi chừng ba mươi cây số.

Tôi dẫn bốn cô gái Thế hệ thứ Tư xuống xe buýt rồi đi bộ vào con đường dẫn tới kho. Thời tiết đầu hè bắt đầu oi ả. Suốt dọc đường đi, cả bốn đứa cứ ríu rít ôn lại những kỷ niệm ngày xưa với Coco, nét mặt ai nấy đều rạng rỡ niềm vui khôn xiết.

Vừa bước tới trước cánh cổng sắt lớn sơn màu xanh lam của nhà kho, một bóng dáng quen thuộc đang đứng vẫy tay gọi lớn: ``\textbf{Aniki! Mọi người ơi!}''

Tôi giật bắn mình. Đó là Kaigou!

``Sao em lại ở đây? Tưởng em đang ở trường mà!?'' tôi trố mắt hỏi.

Kaigou phồng má, khoanh tay trước ngực: ``Thì, hôm nay không có tiết ở trường, mẹ gọi em xuống phụ bốc hàng chứ sao! Anh tưởng mẹ chỉ nhờ vả mỗi anh chắc? Đừng có bảo là anh muốn đuổi em về để một mình chiếm công đấy nhé!''

Tôi thở dài bất lực trước sự lầy lội muôn thuở của cô em gái sinh đôi mỗi khi chỉ có hai chúng tôi. Cả đám kéo nhau vào trong kho. Mẹ tôi đang một mình kéo những tấm vách thạch cao nặng trịch từ thùng xe tải xuống.

Thấy chúng tôi đến, mẹ tôi tươi cười xoa đầu tôi rồi chào đón các cô gái: ``Kuon đến rồi đấy à? Mấy cháu gái này cũng tới phụ cô sao? Cảm ơn các cháu nhiều nhé, hôm nay nhiều hàng lắm đấy!''

Nói là làm, cả nhóm xắn tay áo vào việc. Và ngay trong đợt bốc dỡ đầu tiên, Kanata đã khiến toàn bộ công nhân trong kho phải đứng hình kinh ngạc: Một mình em ấy dùng tay không nhấc bổng ba tấm vách thạch cao dày cộp cùng một lúc, thoăn thoắt khiêng vào trong góc kho mà không hề biến sắc. Đây vốn dĩ là việc mà chỉ có Kaigou làm được.

``Đúng là Khỉ đột Kanata có khác-nora!'' Luna trầm trồ vỗ tay.

``Em đâu có gì đặc biệt đâu, chỉ là lực tay hơi khỏe hơn người thường thôi mà!'' Kanata vội vàng thanh minh, má em dần dần ửng hồng lên vì ngượng ngùng khi cả kho đều quay nhìn với ánh mắt kinh ngạc. Vừa dứt lời, em xoay người sang quát lớn vào đám người đang đứng hình kia: ``Cả hai đứng đực ra đó làm gì thế hả?! Còn mau mau lại phụ tôi khiêng nốt mấy thanh sắt trụ còn lại vào góc kho đi, chừng nào xong hết việc mới được nghỉ!''

``Z-Zâng!'' cô nàng công chúa xứ Kẹo vừa vội vàng đáp lời vừa lao đến nhanh như chớp, hai vòng tay ôm chặt lấy chiếc xô nhựa lớn đựng đầy đinh sắt nặng trịch. Đôi chân nhỏ nhắn vốn mảnh mai của em khẽ khuỵu xuống một chút khi cố gắng rướn người để nâng cái xô lên, ``Ơ\dots\ nặng quá đi-nora\dots tui hơm ngờ nó nặng đến thế này!''

Tôi đứng bên ngoài, nhìn nàng Thiên sứ đang cầm tay chỉ việc, không kìm được mà nói: ``Không chỉ khỏe mạnh đâu, em ấy còn là một vị chỉ huy trưởng tài ba nữa chứ, vừa vào là có thể lập tức điều phối công việc cho cả nhóm.

Kaigou:	Khiếp thật, đúng là khiếp thật\dots\ Chẳng phải đã có mình em ở cái kho này đủ dữ dằn rồi sao, nay lại thêm Kanata-san, thế là có tới hai bà “Sư tử Hà Đông” trong nhà kho rồi còn đâu nữa.

Tôi nháy mắt tinh nghịch, nghiêng đầu cười khúc khích nói với em gái: ``Đâu chỉ có hai đâu, ba bà Sư tử Hà Đông luôn rồi đấy. Cái kho này chắc chắn sẽ sôi nổi lắm đây.''

Kanata vừa đặt nhẹ mấy thanh sắt cuối cùng xuống nền đất, vừa giơ tay chỉ thẳng vào bản thân với vẻ mặt đầy thắc mắc và hơi bực bội: ``Này này! Hai anh chị đang đứng đó lẩm bẩm nói lắp về ai thế hả?! Đừng có nói xấu sau lưng người ta chứ!''

\ct

Khi tất cả các lô hàng cuối cùng cũng đã được xếp vào đúng vị trí, tôi lau vội mồ hôi trán rồi bước đến bên mẹ, giọng thở hổn hển: ``Sắp xếp xong hết rồi mẹ ạ. Còn có chuyến xe nào cần xếp nữa không ạ?''

Bà cười tươi, vỗ nhẹ vào vai tôi: ``Không còn chuyến nào nữa đâu con. Mấy đứa cứ vào văn phòng gần đây nghỉ ngơi đi cho mát.'' Nghe vậy, bốn cô gái Gen 4 liền túa vào căn phòng văn phòng nhỏ để xả hơi sau buổi làm việc vất vả.

Đúng lúc tôi cũng định bước theo sau, giọng mẹ vang lên từ phía sau: ``Kuon, lại đây mẹ bảo cái này.''

Tôi dừng chân quay lại, thắc mắc hỏi: "Dạ, có chuyện gì thế mẹ?"

Không đáp lời, mẹ liền rút ra một cuộn khăn giấy từ trong túi áo làm việc, kéo tôi lại gần và bắt đầu lau sạch vết bụi bặm trên mặt: trán, má, cổ rồi đến cả gáy cũng không bỏ sót. Tôi ngượng chín mặt, khẽ lảng tránh một chút: ``Mẹ ơi, con đâu còn là đứa trẻ con năm ba tuổi nữa đâu mà phải lau mặt như vậy\dots''

Dù nói vậy, tôi vẫn đứng yên để mẹ lau thoải mái, không hề có chút khó chịu nào. Phía bên cạnh, Kaigou tựa lưng vào cửa đang tủm tỉm cười khúc khích nhìn cảnh tượng này.

Mẹ tôi thở dài giả vờ, giọng ôn tồn: ``Cả ngày vất vả thế này chắc mệt lắm nhỉ? Trời ơi sao con tôi lại khổ thế không biết\dots''

Tôi vội vã xua tay, má vẫn còn hơi ửng đỏ: ``C-Con quen rồi mẹ ơi! Mẹ đừng lo cho con, mấy việc này có gì vất vả đâu!''

Kaigou: bật cười lớn, khoanh tay trước ngực trêu chọc: ``Ôi trời ơi, tình cảm mẹ con của nhà mình thật thiêng liêng và cảm động quá đi mất\~{}''

Tôi thở dài ngao ngán, lắc đầu nhìn cô em gái: ``Cái sự thiêng liêng đấy mà tôi thấy phiền phức không chịu được. Đi thôi, vào trong kia nghỉ ngơi cho mát.''

Sau khi mẹ lau xong hết vết bụi bẩn trên người, tôi và Kaigou mới cùng nhau đi vào phòng cùng với mọi người.

Vừa bước qua ngưỡng cửa văn phòng, một tràng quát tháo đầy uy lực bằng chất giọng lanh lảnh vang dội khắp gian phòng: ``\textbf{Cậu kia! Đừng có ngồi lì đó bấm điện thoại nữa! Mau ra kiểm tra lại mã vận đơn đi!}''

``V-Vâng, thưa cô chủ!'' anh thanh niên giật mình, vội chạy về phía xe chở hàng đang được bốc hàng xuống.

``\textbf{Cả cậu lái xe nâng đằng kia nữa! Đưa xe vào góc kệ số ba xếp pallet ngay cho tôi!}''

``Vâng, thưa sếp!'' một anh thanh niên khác đáp lại, cẩn thận đưa chiếc xe lùi về đúng vị trí.

Đứng giữa văn phòng là một người phụ nữ cao ráo với phong thái cực ngầu: Mái tóc ngắn màu tím khói buộc đuôi ngựa gọn gàng, đeo kính gọng đen, diện áo croptop trắng ôm sát để lộ vòng eo săn chắc cùng chiếc quần dài ống rộng. Trên tay cô lăm lăm một chiếc gậy bóng chày có gắn đinh sắt giả.

``Thôi nào mọi người, đừng có đứng đực ra đó nữa, bình tĩnh mà tiếp tục công việc đi chứ\dots'' mẹ tôi cũng lên tiếng thúc giục, rồi quay sang cô nàng, ``Kson-chan, cháu dọa các bạn nhân viên đến thế kia là sao, làm họ sợ hết hồn cả rồi này!''

``Cháu làm việc ở đây bao lâu rồi mà, tính cách nó thế rồi cô ạ\dots\ không làm vậy thì mấy cậu ấy lại trễ tiến độ mất\dots'' cô gái cười khì khì, hai tay vẫn chắp sau lưng, cây gậy bóng chày lắc lư nhẹ theo nhịp chân.

Đó chính là Kson -- người phụ nữ đã chọn một diện mạo, công việc và cái tên mới cho đời mình. Bốn thành viên Gen 4 không thể biết điều đó chỉ bằng một cái nhìn; họ chỉ cảm thấy trong chất giọng, trong dáng đứng ngang tàng và cả câu gọi ``Kuon-paisen'' có một điều gì quen thuộc đến nhói lòng.

Cuộc gặp hôm nay không phải để kéo một người đã rời đi trở lại. Nó là cơ hội để các em gặp Kson, rồi tự quyết định liệu tình bạn cũ có đủ rộng để ôm lấy con người hiện tại của cô hay không.

\pov{Amane Kanata}

Chúng tôi đứng sững sờ ngay ngưỡng cửa. Nếu tôi không tìm hiểu kĩ từ trước, có lẽ tôi đã đấm anh Kuon một trận và quyết định đi về. Lúc đó có thể tôi đã rất tức giận vì đã ăn phải một quả bịp.

Trước mặt chúng tôi là ``Kson'' - người phụ nữ mà Nene đã miêu tả.

Vóc dáng cao lớn ấy. Từng cử chỉ ngang tàng đậm chất giang hồ ấy. Và đặc biệt là chất giọng mang âm sắc rồng quen thuộc không lẫn vào đâu được!

Thế nhưng\dots\ diện mạo bên ngoài của cô ấy lại hoàn toàn là một con người xa lạ. Không có sừng, không có đuôi rồng.

Towa lập tức chồm người về phía trước hai bước, hai mắt tròn xoe ngước nhìn Kson từ đầu đến chân, miệng thì thầm với giọng còn chưa hết bàng hoàng: ``Vậy\dots\ đây chính là Kson-san mà Nene-chan đã kể sao\dots?''

Luna vội vàng trườn ra sau lưng tôi, hai bàn tay nhỏ nắm chặt lấy vạt áo thiên sứ của tôi, giọng nói run rẩy nhưng vẫn giữ được thói quen đuôi câu: ``Trông cô í đáng xợ thựt xự quá-nanora\dots\ Cầm cây gậy to đùng kia trông như sắp đựp ai bất cứ lúc nào í!''

Watame cũng tụt lại sát cạnh cô ấy, cả hai cùng trốn sau bóng lưng của tôi và Towa, nàng Cừu chỉ dám nhô nửa đầu ra hỏi mẹ của anh Kuon với khuôn mặt hết sức lo lắng: ``Kaien-san ơi\dots\ Thế công ty vật liệu này của cô có phải được xã hội đen bảo kê đúng không ạ\dots?''

Bà Kaien bật cười lớn thành tiếng, hai bàn tay vẫy lia lịa về phía mấy cô bé để trấn an, đầu lắc qua lắc lại liên tục: ``Xã hội đen gì đâu cháu ơi, cháu nghĩ quá rồi! Kson-chan chỉ là quản lý kho phụ cô thôi mà, tính cách nó hơi dữ dằn, hay la mắng nhân viên thôi chứ không phải giang hồ đâu mà sợ!''

Kaigou đứng bên cạnh thì liền kéo nhẹ vạt áo làm việc của mẹ, vai còn run lên cười khúc khích trước phản ứng đáng yêu của bốn cô gái, miệng thì phụ họa thêm: ``Mẹ đừng có nói vậy chứ, nhìn Kson-chan đang cầm cây gậy bóng chày kia kìa, phong cách đó đúng chuẩn `ân xạ' xã hội đen rồi còn đâu. Người nào mới đến lần đầu tiên mà không giật mình, không sợ hãi mới là lạ đó mẹ!''

Nghe thấy tiếng nói chuyện, Kson quay đầu lại nhìn chúng tôi, nheo mắt hỏi với vẻ mặt tỉnh bơ: ``Mấy cô gái này là ai đây, Kuon-paisen?''

Cả bốn đứa chúng tôi đồng loạt trố mắt.

``Kuon-paisen.'' Chỉ có một người duy nhất trên đời gọi anh Kuon bằng cái danh xưng kỳ quặc đó!

``Bà đang giả vờ đấy à, Coco-chan?!'' tôi tiến lên một bước, gặng hỏi.

Kson nhún vai, nụ cười tinh nghịch thoáng qua rồi dịu xuống: ``Coco đã khép lại rồi, Kanatan. Tôi là Kson. Nhưng đúng, tôi biết vì sao bà lại hỏi như vậy.''

Towa nhăn mặt: ``Đừng có xạo! Giọng nói với cách gọi Kuon-paisen của bà rành rành ra đó kìa!''

Watame cũng phụ họa: ``Đúng đó! Bọn tôi nhớ cậu gần một năm trời chưa đủ sao?!''

Kson nhìn lần lượt bốn người. Cô vẫn cố nhếch môi như muốn biến mọi chuyện thành một trò đùa, nhưng bàn tay cầm gậy bóng chày đã thả lỏng. Kaigou đứng cạnh chỉ khẽ nói: ``Đừng ép Kson phải trả lời theo cách người khác muốn. Cứ để cô ấy nói điều cô ấy sẵn sàng nói.''

Tôi hít một hơi thật sâu, ra hiệu cho ba người kia đừng chen vào. ``Vậy thì bọn tôi không cần bắt bà phải là Coco. Nhưng xin nói cho bọn tôi biết một điều thôi: Kson có còn muốn gặp bốn đứa bọn tôi không?''

Dứt lời, tôi nháy mắt với anh Kuon và chị Kaigou rồi quay ngoắt người, sải bước dài về phía cửa kho. Towa, Watame và Luna dù chưa hiểu chuyện gì cũng vội vàng bước theo tôi.

Kson khựng lại. Rồi một bàn tay thon dài nhưng rắn rỏi bất ngờ chụp chặt lấy bờ vai tôi từ phía sau. Giọng nói trầm ấm ấy mang theo tiếng cười trừ đầy bối rối: ``Đi sớm thế làm gì\dots\ Kanatan?''

Tôi từ từ quay đầu lại.

Kson đứng đó, gãi gãi sống mũi, nụ cười tinh quái quen thuộc nở rộ trên môi.

Luna rướn cổ lên nhìn kỹ từ đầu đến chân, đôi mắt to tròn bừng sáng rồi lại dè dặt hỏi: ``\dots Kson-san có thật sự là người mà bọn em từng biết không-nora?''

Cô nàng tóc tím lập tức reo lên vui sướng, tay vỗ mạnh vào đùi: ``Ồ! Em bé Công chúa xứ Kẹo cũng có mặt à!''

Watame từ sau lưng tôi ló đầu ra, nước mắt còn rơm rớm trên má: ``Bọn tôi nhận ra những điều quen thuộc, nhưng bọn tôi cũng muốn biết Kson-san muốn được bọn tôi gọi là ai.''

Kson quay sang nhìn nàng Cừu, cười tinh nghịch chỉ tay: ``Cả đùi cừu non Watame cũng đến đông đủ rồi đây!''

Tôi và Kaigou cùng bật cười lớn trước câu đùa cũ quen thuộc, trong khi Watame lập tức nhảy lên không, hai tay chống hàm bực bội: ``Tôi không phải là đồ ăn đâu mà cứ đòi gọi đùi cừu mãi thế!''

Cô nàng chưa dừng lại, lại quay sang nhìn người cuối cùng, hô to vui vẻ: ``Và còn cô nàng TMT luôn đây!''

Towa ngay lập tức xòe tay bịt miệng cô ấy, nước mắt đã trào ra từ lúc nào không hay, vừa sụt sịt vừa la hét: ``Đừng có mà gọi tôi bằng cái biệt danh kỳ quặc đó nữa chứ! Bao lâu rồi mà cậu vẫn không bỏ được thói xấu đó hả?!''

Kson buông cây gậy bóng chày xuống nền nhà, thở một hơi dài chứa đựng bao nỗi nhớ thương. ``Tôi đã chọn rời khỏi cái tên Coco. Tôi không muốn bị kéo ngược về sống trong con người cũ nữa.'' Cô ngước lên, nụ cười trên mặt mềm ra hết mức. ``Nhưng những tháng ngày với mọi người là thật. Mấy câu đùa ngốc nghếch này cũng là thật. Nếu các bà còn muốn có Kson trong Gen 4, thì tôi ở đây.''

``\textbf{Good afternoon, motherfuckers\~{}!}'' cô nàng giang rộng hai cánh tay, hô to câu chào cũ bằng một giọng đầy trêu chọc.

Anh Kuon và chị Kaigou lập tức đồng thanh bóc mẽ: ``Buổi chiều rồi, phải chào `Good afternoon' mới đúng chứ! Vói cả, chú ý lời ăn tiếng nói chút đi.''

Cả văn phòng phá lên cười. Không phải vì Coco đã trở lại, mà vì Kson vẫn giữ một góc tinh nghịch mà cả bốn người đều nhớ.

Và không để cho Kson kịp nói thêm lời nào nữa, Towa, Watame và Luna đã đồng loạt lao tới như những mũi tên, ôm chầm lấy cô. Ba cô bạn vùi mặt vào ngực người bạn cũ mà khóc nức nở: ``Bọn tôi nhớ bà lắm đấy!'' ``Đồ tồi nhà bà!'' ``Lần sau đừng biến mất mà không nói gì nữa!''

Kson vòng hai cánh tay ôm ghì lấy ba người, vừa cười vừa vỗ về lưng từng đứa, khóe mắt cô cũng đã ươn ướt lệ.

Chỉ còn tôi đứng lại cách đó vài bước chân. Là nhóm trưởng của Gen 4, tôi cố giữ cho mình vẻ điềm tĩnh của một người trưởng thành. Tôi khoanh tay, hất cằm nhìn cậu ấy: ``Thế nào, Thằn lằn bự con?''

Kson buông ba người kia ra, bước tới trước mặt tôi, cúi đầu nhìn xoáy vào mắt tôi: ``Kanatan, bà không định ôm tôi một cái à?''

``Thôi, tôi sợ lực tay của tôi làm gãy xương sườn của bà mất,'' tôi bĩu môi.

``Thế mà ngày chia tay trên biển, ai là người ôm chặt cứng lấy tôi khóc nhiều nhất thế nhỉ? Khỉ đột Kanatan nay yếu đuối thế sao?''

``Đừng có gọi tôi là Khỉ đột!''

Hai đứa chúng tôi trừng mắt nhìn nhau chằm chằm vài nhịp, chẳng ai chịu nhường ai, rồi bỗng dưng cùng lúc bật ra tiếng cười lớn, rung cả lồng ngực. Tôi dang rộng hai cánh tay, Kson cũng nghiêng người về phía trước, và chúng tôi siết chặt lấy nhau trong một cái ôm thật chặt. Cái ôm ấy không xóa đi cuộc chia tay năm nào, cũng không buộc cô ấy quay lại làm Coco. Nó chỉ xác nhận rằng, sau bao tháng ngày, Kson vẫn muốn để bốn đứa chúng tôi ở trong đời mình. Hơi ấm từ người bạn cũ lan tỏa, làm cho những giọt nước mắt tôi cố gắng kìm nén bấy lâu nay bỗng dưng trào ra, ướt đẫm vai áo của cô ấy.

Từ góc văn phòng, anh Kuon đứng dựa vào tường, hai tay khoanh trước ngực, khóe miệng nhô lên một nụ cười ấm áp không thể che giấu: ``Họ cuối cùng cũng được đoàn tụ thật sự rồi nhỉ.''

Chị Kaigou dựa vai vào anh, tay chỉ nhẹ nhàng về phía nhóm năm cô gái đang quấn quýt lấy nhau, mắt cũng rạng rỡ niềm vui: ``Nhìn cảnh đó thì chắc chắn là thế rồi. Mấy đứa đã chờ khoảnh khắc này quá lâu.''

Bà Kaien đứng cạnh hai người con, lấy khăn giấy lau nhanh một giọt lệ ươn ướt ở khóe mắt, miệng vẫn mỉm cười hạnh phúc: ``Lần đầu tiên mẹ thấy Kson cười tươi như vậy kể từ ngày nó chuyển về đây làm việc, đúng là có bạn bè bên cạnh thì khác.''

Đứng sau lưng ba người, hai nhân viên trong kho nhìn cảnh tượng đầy nước mắt và nụ cười trước mặt hoàn toàn ngẩn ngơ, chẳng hiểu chuyện gì đang diễn ra.

``Ờm, thưa sếp\dots\ vừa rồi có chuyện gì xảy ra vậy ạ? Tôi\dots\ tôi hơi không theo kịp.''

``Vâng ạ, chúng tôi\dots\ chúng tôi cần một lời giải thích. Sao tự nhiên mọi người lại khóc rồi ôm nhau thế ạ?''

Bà Kaien lập tức lau sạch vết lệ trên mặt, lấy lại vẻ nghiêm túc thường ngày, quạt nhẹ tay đuổi hai người: ``Đi làm việc đi, việc của công ty không cần các anh phải tò mò. Còn chuyến hàng tới \ruby{Hokunei}{Bắc Ninh} còn đang chờ xếp lên kệ đấy, còn chần chừ gì nữa!''

``V-Vâng thưa bà! Chúng tôi đi làm bây giờ!''

\dots Hãy gác lại những câu hỏi không đáng có của những người ngoài cuộc, hãy để khoảnh khắc riêng này thuộc về những người đã chờ đợi nó suốt gần một năm qua.

Trong vòng ôm ấm áp của Coco, tôi - Amane Kanata, nhóm trưởng của Thế hệ thứ tư, cảm nhận từng ngọn lửa ấm lan tỏa trong tim. Khoảnh khắc này, tôi biết rằng gia đình holoForce của chúng tôi cuối cùng cũng đã trọn vẹn trở lại. Năm người một nhà, không còn thiếu ai nữa.

Coco à\dots\ tôi đã mòn mỏi chờ đợi ngày này từ cái ngày chúng ta đứng trên bờ biển chia tay, từ cái ngày em bước ra khỏi căn hộ chung của bọn tôi. Tôi nhớ em từng đêm, nhớ những câu chuyện đùa cợt, nhớ những lúc chúng ta cùng nhau vượt qua sóng gió. Cái con thằn lằn bự con hay trêu chọc tôi này, tôi đã nhớ nó đến phát điên.

Tôi thầm cầu nguyện rằng từ nay, chúng ta sẽ không còn phải xa nhau nữa, sẽ lại là những người bạn bình thường vui vẻ như ngày xưa, sẽ lại cùng nhau stream, cùng nhau cười đùa, sẽ không còn những khoảng cách vô hình ngăn cách chúng ta như trong bao tháng ngày qua.

Mừng bà đã trở về, Kiryu Coco - hay Kson, người quản lý kho đầy ngầu của chúng tôi. Bọn tôi đã đợi bà quá lâu rồi.

%==========================================TRUYỆN PHỤ=========================================
\omk

\pov{Hikawa Kuon}

Sau màn hội ngộ đẫm nước mắt và nụ cười, cả cái nhà kho bỗng chốc biến thành chiến trường dưới sự chỉ huy của ba ``nữ tướng'': mẹ tôi, Kanata và Kson.

Mẹ tôi phụ trách điều phối xe tải, và việc điều hành công việc này cũng chẳng có gì khó với tư cách là một chủ kho.

Mẹ tôi vừa gõ gõ đầu bút lên màn hình máy tính bảng, vừa vẫy gọi nhân viên phụ bốc hàng đang đi ngang qua: ``A này, lại đây giúp chị việc này cái\dots\ Em có thể xếp lịch chuyển đơn hàng này sang quận Ryuuhei được không?''

Nhân viên nhanh chân bước lại, cúi đầu xem màn hình mẹ tôi chỉ: ``Dạ được ạ. Khoảng bao nhiêu hàng hóa cần vận chuyển chị?''

``Số lượng vách thạch cao và xà gồ (tôi vẫn gọi theo thói quen cũ là ``xương cá'') chị ghi hết chi tiết trong file đính kèm rồi đấy, em kiểm tra lại cho kỹ nhé.''

``Vâng ạ, em ghi nhận rồi sẽ xử lý ngay.''

Mẹ tôi vừa gửi file cho nhân viên, vừa ngoái đầu gọi lớn vào giữa kho: ``Kaigou ơi, lại đây bốc cho mẹ ba tấm vách ở chồng hàng đằng kia\dots\ Mà gọi cả Kanata-chan sang phụ luôn, xem lực của con bé có khỏe như mẹ nghe đồn không nhá!''

Kaigou vừa lau mồ hôi trán, vừa đáp liền từ xa: ``Vâng ạ, con tới đây liền\~{}''

Rồi mẹ tôi quay sang tôi, đang đứng gần đó tán gẫu với Kson, vỗ vỗ vai tôi: ``Kuon ơi, lại đây con. Con ngồi đây làm sổ sách, đối chiếu hóa đơn giúp mẹ nhé.''

Tôi lập tức ngồi vào ghế: ``Ờm\dots\ vâng. Cái việc này đúng là con làm giỏi nhất mà! Mà con phải xử lý bao nhiêu đơn hàng ạ?''

Mẹ tôi đẩy cả một chồng phiếu xuất kho sang trước mặt tôi: ``Chừng khoảng 5, 6 đơn gì đấy, đều là đơn gọn thôi mà.''

Tôi há hốc miệng nhìn chồng giấy dày cộp, thốt lên không thể kìm được: ``Vãi đạn, nhiều thế đâu phải `gọn' đâu mẹ ơi\dots''

Trong khi đó, nàng Thiên sứ Amane Kanata và người quản lý kho đầy cá tính Kson đã nhanh chóng phân chia vai trò, cùng nhau đốc thúc từng khâu bốc dỡ hàng hóa một cách nhịp nhàng. Nhìn từ xa, hai người chẳng khác gì hai bà chủ doanh nghiệp đang ra lệnh cho toàn bộ nhân viên dưới quyền phải chạy đôn chạy đáo, khiến ai cũng phải ngán ngẩm trước tác phong làm việc dứt khoát của họ.

Đặc biệt là Kanata, em ấy dường như đang bộc lộ sự hứng thú bất ngờ với công việc điều phối này, mái tóc ngắn tung tăng mỗi khi bước đi vội vàng, mắt lướt qua từng góc kho để đảm bảo không có khâu nào bị trễ.

Kson vỗ mạnh vào đống thanh sắt xếp chồng ở góc kho, cất giọng hô to về phía Kanata đang kiểm tra số lượng hàng trên kệ: ``Kanatan! Qua đây phụ tôi chất mấy thanh sắt này lên xe tải cái!''

Kanata lập tức xoay người, tay vỗ vào đùi đầy hào hứng: ``Rồi đây!'' Ngay sau đó em chỉ về phía Luna đang đứng nép một bên, giọng đầy ra lệnh: ``Luna, điều khiển xe nâng đến kệ hàng số 3 ở đằng kia, lấy đủ 8 thanh sắt cùng loại mang đến đây!''

Luna giơ tay chào kiểu quân đội, miệng reo lên đầy nhiệt tình: ``Rõ thưa chỉ huy-nanora!'' Em nhảy tót lên ghế lái xe nâng, tay lướt qua các cần điều khiển một cách hào hứng.

Kson quay sang một nhân viên đang đi ngang qua, tay vẫy gọi mạnh mẽ: ``Cậu kia, dừng lại một chút! Lập tức chuẩn bị đơn hàng số 78 của tôi, đừng để bị sót mặt hàng nào!''

Cậu nhân viên giật mình, vội vàng đáp lại vừa chạy đến chỗ Kson: ``V-vâng thưa sếp! Em làm ngay!''

Một nhân viên khác thấy vậy liền nhanh chân chạy theo, tay giơ lên đề nghị: ``Để em phụ anh ấy kiểm tra mã vận đơn cho nhanh ạ!''

Kson gật đầu hài lòng, tay vỗ vào vai hai người: ``Đúng mới là làm ăn! Làm tốt đi, xong việc tôi trà đá mời cả đám!''

Bên cạnh đó, Kanata lại gọi về phía Watame đang đứng ngẩn ngơ xem Luna điều khiển xe nâng: ``Watame! Bà ra kia mà lùi gà vào chuồng đi! Đừng đứng đó làm cảnh nữa!''

Watame quay người, mặt đầy ngơ ngác hỏi lại: ``Gà nào? Ở đây có nuôi gà à Kanata-chan? Sao tôi không thấy đâu cả?''

Luna vừa xoay vô lăng xe nâng để điều chỉnh vị trí, vừa cười khúc khích giải thích giúp nàng cừu non: ``Ý bà í là dẫn hứng cho xe tải lùi vào zị trí xếp hàng đó-nora! Lùi xe hay nói dân gian là lùi gà vào chuồng mà!''

Nàng Cừu mặt bừng lên vì ngượng, tay vò mái tóc vàng bồng bềnh của mình: ``R-Rồi! Làm tôi tưởng mình phải đi dọn chuồng gà thật, còn phải hót phân chiêm chiếp nữa chứ\dots'' Em vội vàng chạy ra phía sau xe tải, tay vẫy lia lịa để hướng dẫn tài xế.

Kson lại tiếp tục quét mắt khắp kho, dừng lại ở Towa đang đứng bấm điện thoại nhân viên, liền gọi lớn: ``Towa! Còn đứng đó làm gì?!''

Towa giật mình, điện thoại suýt rơi khỏi tay, miệng lắp bắp đáp: ``S-Sao vậy Kson? Có việc gì cần tôi làm ạ?''

``Vào trong văn phòng lấy đống đơn hàng đã đối soát từ Kuon-paisen đi, xong xuôi thì mang tất cả ra đây để tôi ký nhận! Đừng có đi đường vòng mà đánh rơi giấy tờ đâu đấy!'' cô nàng Tổng trưởng nhắc nhở, tay chỉ về phía căn phòng văn phòng ở cuối kho.

Towa vội vàng cất điện thoại vào túi quần, bước nhanh vào trong: ``R-Rồi, tôi làm ngay đây!''

Tôi bị đày vào một góc bàn làm việc để xử lý cả một chồng hóa đơn và phiếu xuất kho dày cộp. Ngồi cạnh tôi, Kaigou đang loay hoay với chiếc máy tính bảng, bỗng huých vai tôi: ``\vie{Anh Kuon này\dots\ Tính giá trị trung bình trên trang tính dùng hàm gì ấy nhỉ?}''

Tôi nhíu mày nhìn sang: ``\vie{Anh tưởng em là dân IT chuyên nghiệp cơ mà? Mấy cái cỏn con này mà phải hỏi à?}''

``\vie{Em chuyên viết code hệ thống chứ có chuyên làm bảng tính Excel đâu!}'' Kaigou phụng phịu.

``\vie{Lên mạng mà tra. Anh còn cả đống đơn hàng phải đối soát đây này,}'' tôi gắt lên. Kaigou bĩu môi lườm tôi một cái cháy mặt, lẩm bẩm: ``\vie{Mừ\~{} Đồ xấu tính\dots}''.

Trong khi đó, ở ngoài bãi hàng, những tình huống dở khóc dở cười liên tục diễn ra.

Luna xung phong nhận nhiệm vụ lái xe nâng điện để chuyển các thanh sắt trụ. Cô bé kéo mũ bảo hộ xuống, hí hửng bước lên ghế lái: ``Ủa, cái cần gạt này kéo xuống là hạ càng nâng đúng không-nora?''

Kson đứng gần đó đáp: ``Ừ, kéo nhẹ xuống là được!''

*CẠCH* ``Như thế nì ạ?''

*VÈO! RẦM!*

Luna kéo nhầm cần số lùi và nhấn ga quá đà. Chiếc xe nâng vọt lùi về phía sau với tốc độ tên lửa, húc thủng một mảng tường thạch cao lớn rồi lật nghiêng sang một bên.

Mẹ tôi và Kaigou hốt hoảng chạy ra đỡ cô bé dậy, may mắn là nàng công chúa kẹo ngọt không hề xây xát gì ngoài việc mặt mày lấm lem bụi vôi.

Ở một góc khác, Watame nhận nhiệm vụ làm lơ xe điều hướng cho một chiếc xe tải chở hàng lùi vào cửa kho. ``Lùi sang bên này chú ơi!'' em ấy đứng sau xe huơ tay lia lịa.

Tài xế thò đầu ra ngoài quát: ``\textbf{Bên nào?! Bên trái của chú hay bên trái của cháu?!}''

Nàng Cừu cuống quá hét toáng lên: ``\textbf{Bên trái của cháu ạ!}''

Chiếc xe tải lập tức đánh lái sang phải theo quán tính, và *RẦM* một tiếng, đuôi xe quẹt thẳng vào dãy vách tôn ngăn cách mà Kson đã mất nguyên cả ngày hôm qua để dựng lên. Toàn bộ dãy tôn đổ sập hàng loạt như quân cờ domino.

Kson ôm đầu gào thét: ``\textbf{Ối giời đất ơi! Bức tường ngăn của tôi! Đứa nào chỉ đường đấy?!}''

Watame lập tức chắp hai tay trước ngực, giả nai với khuôn mặt vô tội nhất trần đời: ``Watame hơm biết gì nha\~{} Watame hơm có nhỗi\~{}''

Tài xế nhảy xuống xe cự cãi: ``Rõ ràng là cô bé này xi-nhan cho tôi mà!''

Chứng kiến cảnh tượng hỗn loạn ấy, ba mẹ con tôi cùng Kson nhìn nhau cười trừ. Tôi bước ra vẫy tay gọi cả bốn cô gái Gen 4 lại: ``Mấy đứa có muốn học bí kíp de xe vào chuồng chuẩn chỉ của nhà Hikawa không?''

``Học kiểu gì ạ?'' Cả bốn đứa tò mò hỏi.

Chúng tôi không nói một lời. Tôi, Kaigou, mẹ tôi và Kson bước ra đứng thành một hàng ngang sau đuôi chiếc xe tải tiếp theo.

Và rồi, cả bốn người chúng tôi đồng loạt thực hiện những cử chỉ tay nhịp nhàng, uốn lượn bước lùi theo từng nhịp xe lăn bánh - trông hệt như một điệu nhảy đồng diễn ``vũ điệu de xe'' cực kỳ điêu luyện trên nền nhạc không lời. Thành thật, trông chúng tôi chẳng khác gì những tên ngốc cả.

Bốn cô gái Gen 4 trố mắt kinh ngạc, rồi thích thú bắt chước làm theo y chang. Thế là giữa buổi chiều oi ả trong nhà kho, tám con người cùng nhau múa điệu de xe đầy kỳ quặc nhưng lại chuẩn xác đến từng centimet.

\ct

Ánh vàng hoàng hôn nhuốm đầy cả mái nhà kho, kéo dài những bóng cây trên mặt đất lấm lem bụi đường. Toàn bộ các chuyến hàng trong ngày cuối cùng cũng được xếp gọn vào đúng kệ, lên đúng xe, cứu nguy cho tiến độ mà mẹ tôi đã lo lắng suốt cả tuần một cách ngoạn mục. Mọi người vừa lau mồ hôi trán, vừa tập trung lại ở cửa kho để chia tay trước khi mỗi người về nhà nghỉ ngơi.

Kson chống hai tay vào hông, giọng nói vang vọng sau khi vừa tiễn chiếc xe tải cuối cùng lăn bánh ra cổng: ``Mấy đứa về nhà dưỡng sức đi nhé! Cả ngày lao động cực lắm rồi!''

Bốn người còn lại cùng đáp liền, giọng nói vẫn tràn đầy năng lượng dù chân tay đã mỏi nhừ: ``Rồi! Cũng nhờ có bà chỉ dẫn nữa chứ!''

Nàng Tổng trưởng giơ tay chỉ thẳng vào Kanata, miệng nở nụ cười tinh quái như thời còn chung phòng: ``Nhất là Kanatan đấy! Về nhà ăn nhiều vào, tập thể dục đều đặn, cho ngực to ra một chút! Đẹp gái hơn!''

Thiên sứ lập tức nhảy lên không, hai má bừng lên vì ngượng, tay chỉ thẳng vào mặt cô bạn rồng: ``Này! Này! Này! Bà vừa nói gì? Ý bà là bà chê ngực tôi lép hả?! Sao dám nói chuyện đó trước mặt mọi người thế!''

Kson ôm bụng cười lớn, tiếng cười vang cả khoảng sân kho: ``Úi, bại lộ thật rồi! Tôi tính giấu nhẹm đi mà\dots\ hà hà!''

``Cái bà này! Dám trêu chọc nhóm trưởng à?! Không tha cho bà đâu!'' Kanata la lên, vừa chạy tới vừa định vồ lấy áo của cô bạn.

Ngay lập tức, cô nàng quản lý kho xạ xã hội Kson liền quay người bỏ chạy, còn Kanata thì đuổi theo sát gót, hai người vòng qua vòng lại các dãy kệ hàng còn trống, tiếng cười nói đùa giòn tan của hai đứa va vào các bức tường, vang vọng khắp cả gian nhà kho rộng lớn. Cô vừa chạy vừa ngoái đầu trêu chọc thêm, còn Kanata thì vừa thở hổn hển vừa cố gắng bắt kịp người bạn cao lớn của mình.

Towa, Watame và Luna đứng sát bên tôi, ba người chống tay vào nhau, ánh mắt lấp lánh ngập tràn niềm hạnh phúc khi nhìn cảnh tượng quen thuộc ấy. Từ cái ngày họ phải chia tay nhau ở bờ biển, họ chưa bao giờ thấy Kson và Kanata có thể đùa giỡn vô tư như thế này nữa. Họ không cần gọi sự ngang tàng của Kson là một chiếc mặt nạ; chính trong con người hiện tại ấy, những ký ức chung vẫn có chỗ để sống.

Ác ma vắt tay lên ngực, khóe mắt hơi ươn ướt nhưng miệng thì nở nụ cười ấm áp: ``Họ đúng là đã quay trở về như trước nhỉ. Chẳng có gì thay đổi cả.''

Công chúa xứ Kẹo chắp hai tay trước bụng, đầu gật gù liên tục, giọng nói vẫn giữ được thói quen đuôi câu quen thuộc: ``Phải thành thực là họ thân nhau thựt đấy-nora\dots\ Tôi vừa mừng vừa muốn khóc quá.''

Nàng Cừu lau nhanh giọt lệ đang tràn ra ở khóe mắt, gật đầu lia lịa đồng tình: ``Đúng đó! Kson vẫn là bạn của bọn mình. Không có gì thay đổi được tình bạn của bọn mình!''

Tôi tựa lưng vào khung cửa kho, ngắm nhìn năm cô gái của Thế hệ thứ Tư.

Họ thực sự đã trưởng thành. Dù không còn đứng chung dưới một mái nhà công ty, nhưng ngọn lửa tình bạn giữa họ vẫn luôn cháy rực, bền chặt và vĩnh cửu.

Tôi thầm mỉm cười nhẹ nhõm. Chiến dịch Réunion dành cho Gen 5 và Gen 4 đã thành công mỹ mãn ngoài mong đợi.

Thế nhưng, khi màn đêm buông xuống, một nỗi trăn trở mới lại bắt đầu len lỏi vào tâm trí tôi.

\textit{holoFantasy} - thế hệ thứ Ba.

Tôi ngước mắt nhìn về phía xa xăm của thành phố.

Câu chuyện của Mikeneko, người đang ẩn mình trên tầng tám của căn nhà chung của gia đình tôi sẽ là một thách thức khó nhằn hơn gấp trăm lần tất cả những gì tôi vừa trải qua. Những vết thương tinh thần trong lòng cô ấy vẫn còn mới nguyên, niềm tin vào thế giới xung quanh đã bị bẻ gãy và nứt nẻ đến mức không thể nhận ra. Làm sao tôi có thể vá lại một trái tim đã tan thành từng mảnh nhỏ mà không làm nó càng rạn nứt thêm trong quá trình cố gắng đó?

Ngay cả việc mở lời trò chuyện với một người lạ ngoài vòng tròn gia đình Hikawa đã là điều gần như bất khả thi với cô ấy, huống chi là đối mặt lại với những người đồng đội cũ trong thế hệ \textit{holoFantasy}. Từ khi bước vào căn nhà này, Mike-chan gần như tự giam mình trong phòng, hoàn toàn cắt đứt liên lạc với mọi người bên ngoài bức tường của căn phòng nhỏ đó.

Có lẽ với ca này, tôi phải thực sự đặt hết tâm trí và cẩn trọng từng bước một. Tôi sẽ phải tìm cách thuyết phục mọi người, đặc biệt là những đồng đội cũ của em ấy, rằng Rushia đã thực sự rời xa thế giới này, ít nhất là trong giai đoạn đầu này, cho đến khi tâm lý của Mike-chan đủ ổn định để có thể đối mặt với sự thật phức tạp phía sau.

Tôi thở một hơi dài, những lo âu cứ cuộn tròn trong ngực, nhưng ánh mắt nhìn về phía màn đêm đang phủ dần lên thành phố không hề nao núng.

Một hành trình dài đằng đẵng, chất chứa đầy những khó khăn và thử thách còn đang chờ đợi tôi ở phía trước, và tôi nhất định phải bước qua nó.

\end{document}
